\documentclass[10pt,a4paper]{amsart}
\usepackage[margin=19mm]{geometry}
\usepackage{amsmath,amssymb,mathtools}
\usepackage[scr=rsfs,cal=euler]{mathalfa}
\usepackage{xfrac}
\usepackage[unicode,hidelinks,bookmarks=false]{hyperref}
\newcommand{\ie}{i.e.\ }
\newcommand{\cf}{cf.\ }
\newcommand{\resp}{respectively\ }
\newcommand{\Subsection}{\subsection}
\providecommand{\nobreakdash}{\nobreak\hbox{-}}
\setlength{\emergencystretch}{2em}
\multlinegap0pt
\usepackage{siunitx}
\usepackage{siunitx}
\usepackage{siunitx}
\usepackage{siunitx}
\usepackage{amssymb}
\usepackage{mathtools}
\usepackage{bm}
\usepackage{xcolor}
\usepackage{tikz}
\usepackage[shortlabels]{enumitem}
\usepackage{siunitx}
\usepackage{float}
\newtheorem{theorem}{Theorem}
\title{When Can Depth Replace Precision?\\ A Resource Theory of Quantized Neural Computation}
\author{Mojtaba Soltanalian}
\date{Source: arXiv:2607.23390v1 (CC BY 4.0)}
\begin{document}
\maketitle

\noindent\emph{Source and licence.} When Can Depth Replace Precision?\\ A Resource Theory of Quantized Neural Computation. Version arXiv:2607.23390v1 (CC BY 4.0), distributed by arXiv under the Creative Commons Attribution 4.0 International licence, http://creativecommons.org/licenses/by/4.0/, as recorded in the arXiv licence metadata for this exact version and shown on the abstract page https://arxiv.org/abs/2607.23390v1. Full licence notice: \texttt{licenses/arxiv-2607.23390-CC-BY-4.0.txt}.\par\smallskip

\noindent\textit{Editorial scope.} Counted unit: the article's theorem eq:fixed-teacher-exact-main (source0.tex lines 1442--1463). The article's complete original proof of it is delivered with it (source0.tex lines 1465--1526, one \texttt{proof} environment of 62 lines). No mathematics is written, repaired or re-derived: the statement and the proof are the article's own lines, in the article's own notation, and no retained line is substituted or re-flowed.  Retained uncounted as the source-local context the proof needs: lines 1430--1434; lines 1436--1439.  Nothing was omitted from any retained span, so the package carries no omission record.  No retained line contains a citation, so the wrapper carries no bibliography and no bibliography entry.  No retained line contains a figure environment, an image-inclusion command or drawing code, so no image is delivered and the package declares no asset dependency.  Editorial formatting only, in addition: an AMS \texttt{amsart} preamble that restates the article's own declarations and macros used by the retained lines, namely the environments \texttt{theorem} on separate counters, together with the standard packages those lines need.  The retained environments print in this document as Theorem 1 in order of appearance, whereas the article prints the same items with the numbering of its own full document.  The complete native source of the article as distributed in the arXiv e-print for this exact version is bundled unchanged as \texttt{sources/206023/source0.tex}.
\begin{equation}
 x_{k+1}=\left(1-\frac1D\right)x_k+\frac{q_k}{D},
 \qquad q_k\in\{0,a\},\qquad x_0=0.
 \label{eq:binary-recursion-main}
\end{equation}

\begin{equation}
 x^\star=a(1-e^{-1})
 \label{eq:fixed-teacher-target-main}
\end{equation}

\begin{theorem}[Exact best error for one fixed binary teacher]
\label[theorem]{thm:fixed-teacher-main}
For every integer $D\ge2$, the all-$a$ schedule is a globally closest depth-$D$ endpoint of \eqref{eq:binary-recursion-main} to the fixed target \eqref{eq:fixed-teacher-target-main}.  The exact optimum is
\begin{equation}
 E_D^{\mathrm{bin}}(x^\star)
 =a\left[e^{-1}-\left(1-\frac1D\right)^D\right].
 \label{eq:fixed-teacher-exact-main}
\end{equation}
Moreover,
\begin{equation}
 \frac{a}{4eD}
 \le E_D^{\mathrm{bin}}(x^\star)
 \le \frac{a}{2e(D-1)},
 \label{eq:fixed-teacher-bounds-main}
\end{equation}
and
\begin{equation}
 E_D^{\mathrm{bin}}(x^\star)
 =\frac{a}{2eD}+O(D^{-2}).
 \label{eq:fixed-teacher-asymptotic-main}
\end{equation}
\end{theorem}

\begin{proof}
Write $r_D=1-D^{-1}$.  Unrolling the recursion gives
\begin{equation}
 x_D=\frac1D\sum_{k=0}^{D-1}r_D^{D-1-k}q_k.
 \label{eq:binary-endpoint-expansion-main}
\end{equation}
The largest attainable endpoint is
\begin{equation}
 M_D=a(1-r_D^D),
 \label{eq:binary-largest-main}
\end{equation}
obtained by choosing $q_k=a$ at every step.  Because the coefficients in \eqref{eq:binary-endpoint-expansion-main} are positive and increase with $k$, the second-largest endpoint is obtained by replacing the earliest, least heavily weighted symbol by zero.  Hence the gap from the largest to the second largest endpoint is
\begin{equation}
 g_D=\frac{a}{D}r_D^{D-1}.
 \label{eq:binary-top-gap-main}
\end{equation}
It is therefore enough to prove that the fixed teacher belongs to the Voronoi cell of $M_D$, namely
\begin{equation}
 M_D-x^\star\le \frac{g_D}{2}.
 \label{eq:binary-voronoi-condition-main}
\end{equation}
Put $u=D^{-1}$.  For $0<u\le1/2$, the convergent series identity
\begin{align}
 1+\left(\frac1u-1\right)\log(1-u)+\log(1-u/2)
 &=\sum_{n=2}^{\infty}\frac{u^n}{n}
   \left(\frac1{n+1}-\frac1{2^n}\right)
 \ge0
 \label{eq:binary-series-main}
\end{align}
uses only $2^n\ge n+1$.  Exponentiating \eqref{eq:binary-series-main} yields
\begin{equation}
 e^{-1}\le r_D^{D-1}\left(1-\frac1{2D}\right).
 \label{eq:binary-exponential-ineq-main}
\end{equation}
Consequently,
\begin{align}
 M_D-x^\star
 &=a(e^{-1}-r_D^D) \\
 &\le \frac{a}{2D}r_D^{D-1}
 =\frac{g_D}{2},
 \label{eq:binary-voronoi-proof-main}
\end{align}
which proves global optimality and the exact formula.

For the quantitative envelope, define
\begin{equation}
 \theta_D=-1-D\log(1-D^{-1})
 =\sum_{n=2}^{\infty}\frac1{nD^{n-1}}.
 \label{eq:theta-main}
\end{equation}
Then
\begin{equation}
 \frac1{2D}\le\theta_D\le\frac1{2(D-1)}\le\frac12
 \label{eq:theta-bounds-main}
\end{equation}
and
\begin{equation}
 e^{-1}-r_D^D=e^{-1}(1-e^{-\theta_D}).
 \label{eq:theta-error-main}
\end{equation}
For $0\le\theta\le1/2$, $\theta/2\le1-e^{-\theta}\le\theta$.  Substitution gives \eqref{eq:fixed-teacher-bounds-main}.  Expanding \eqref{eq:theta-main} and then $1-e^{-\theta_D}$ gives \eqref{eq:fixed-teacher-asymptotic-main}.
\end{proof}

\end{document}
